\begin{center}
	\textbf{Partie A}
\end{center}

Un artisan crée des bonbons au chocolat dont la forme rappelle le profil de la montage locale représentée en \textbf{Figure 1}. La base d'un tel bonbon est modélisée par la surface grisée, définie ci-dessous dans un repère orthonormé d'unité 2~cm (\textbf{Figure 2}).

\begin{center}
\begin{tabularx}{\linewidth}{*{2}{>{\centering \arraybackslash}X}}
\psset{unit=2.cm,arrowsize=2pt 3}
\begin{pspicture}(-1.2,-0.2)(1.4,1.6)

\def\fond{\psplot[plotpoints=2000,linewidth=1.25pt]{0.414214}{1}{1 x dup mul sub 2.71828 x exp mul}}
\rput(0.42,0.3){\fond}
\psplot[plotpoints=2000,linewidth=1.25pt]{-1}{1}{1 x dup mul sub 2.71828 x exp mul}
\rput(0,0.){\pscustom[fillstyle=solid,fillcolor=lightgray]{\psplot[plotpoints=2000,linewidth=1.25pt]{-1}{1}{1 x dup mul sub 2.71828 x exp mul}\psline(1,0)(-1,0)}}
\psline(1,0)(1.45,0.3) \psline(0.35,1.24)(0.85,1.56)
\psline{<->}(1.2,0)(1.65,0.3) \uput[r](1.42,0.1){3 cm}
\end{pspicture}
&
\psset{unit=2.cm,arrowsize=2pt 3}
\begin{pspicture*}(-1.2,-0.5)(1.4,1.4)
\pscustom[fillstyle=solid,fillcolor=lightgray]{\psplot[plotpoints=2000,linewidth=1.25pt]{-1}{1}{1 x dup mul sub 2.71828 x exp mul}\psline(1,0)(-1,0)}
\psplot[plotpoints=2000,linewidth=1.5pt]{-1}{1}{1 x dup mul sub 2.71828 x exp mul}
\uput[ur](0.7,1){$\mathcal{C}_f$}
\psgrid[subgriddiv=5,  gridlabels=0, gridcolor=lightgray] 
\psaxes[linewidth=1.25pt]{->}(0,0)(-1.2,0)(1.4,1.4)
\end{pspicture*}\\
\textbf{Figure 1}&\textbf{Figure 2}\\
\end{tabularx}
\end{center}

Cette surface est délimitée par l'axe des abscisses et la représentation graphique notée $\mathcal{C}_f$ de la fonction $f$ définie sur $[-1~;~1]$ par :
$f(x) = \left(1 - x^2\right)\e^x.$

L'objectif de cette partie est de calculer le volume de chocolat nécessaire à la fabrication d'un bonbon au chocolat.

\begin{enumerate}
\item 
	\begin{enumerate}
		\item 
Pour tout réel $x$, on a $\e^{x}>0$, et si $x\in 	[-1~;~1]$, on a $1-x^2\geqslant 0$. 

Donc pour tout $x$ de $[-1~;~1]$, on a $f(x)\geqslant 0$.
		
		\item 
Soit $I= \ds\int_{-1}^1 f(x) \d x$.

On pose $u(x)=1-x^2$ et $v'(x)=\e^{x}$. Donc $u'(x)=-2x$ et $v(x)=\e^{x}$.

Par intégration par parties: $\ds\int_{a}^{b} u(x)v'(x) \d x = \left [ u(x)v(x)\strut\right ]_{a}^{b} - \ds\int_{a}^{b} u'(x)v(x) \d x$.

Donc $I= \left [ (1-x^2)\e^{x} \strut \right ]_{-1}^{1} - \ds\int_{-1}^{1} (-2x)\e^{x} \d x
= 0+2 \ds\int_{-1}^{1} x\e^{x} \d x =  2\ds\int_{-1}^{1} x\e^{x} \d x$
	\end{enumerate}
	
\item Le volume $\mathcal{V}$ de chocolat, en cm\up{3}, nécessaire à la fabrication d'un bonbon est donné par :
$\mathcal{V} = 3 \times S$
où $S$ est l'aire, en cm\up{2}, de la surface colorée (\textbf{Figure 2}).



On calcule $\ds\int_{-1}^{1} x\e^{x} \d x$ par une intégration par parties.

On pose $u(x)=x$ et $v'(x)=\e^{x}$. Donc $u'(x)=1$ et $v(x)=\e^{x}$.

$\begin{aligned}
\ds\int_{-1}^{1} x\e^{x} \d x
& = \left [ x \e^{x} \strut \right ]_{-1}^{1} - \ds\int_{-1}^{1} 1\times \e^{x} \d x
= \left [ x \e^{x} \strut \right ]_{-1}^{1} - \left [  \e^{x} \strut \right ]_{-1}^{1}
= \left [ 1 \e^{1} - (-1)\e^{-1} \strut \right ] -  \left [ \e^{1} - \e^{-1} \strut \right ]\\
& = \e + \e^{-1} -\e +\e^{-1} = 2 \e^{-1}
\end{aligned}$ 

Donc $S = \ds\int_{-1}^{1} f(x) \d x = 2 \ds\int_{-1}^{1} x\e^{x} \d x = 4 \e^{-1}$.

$\mathcal{V} = 3 \times S = 12\e^{-1}$ donc $\mathcal{V} \approx 4,4$ cm$^3$.
\end{enumerate}

\smallskip

\begin{center}
\textbf{Partie B}
\end{center}

On s'intéresse maintenant au bénéfice réalisé par l'artisan sur la vente de ces bonbons au chocolat en fonction du volume hebdomadaire des ventes.



Ce bénéfice peut être modélisé par la fonction $B$ définie sur l'intervalle $[0,01~;~+\infty[$ par :  $B(q) = 8q^2 [2 -3\,\ln q] - 3.$

Le bénéfice est exprimé en dizaines d'euros et la quantité $q$ en centaines de bonbons.

On admet que la fonction $B$ est dérivable sur $[0,01~;~+\infty[$. On note $B'$ sa fonction dérivée. 

\begin{enumerate}
\item 
	\begin{enumerate}
		\item
\begin{list}{\textbullet}{}
\item $\ds\lim_{q \to +\infty} \ln q = +\infty$ donc $\ds\lim_{q \to +\infty} 2- 3\,\ln q = -\infty$
\item $\ds\lim_{q \to +\infty} q^2 = +\infty$
\item Donc $\ds\lim_{q \to +\infty} 8q^2 [2 -3\,\ln(q)] = -\infty$ et donc $\ds\lim_{q \to +\infty} 8q^2 [2 -3\,\ln(q)]-3 = -\infty$
\end{list}

On a donc démontré que $\ds\lim_{q \to +\infty} B(q)=-\infty$.

		\item Pour tout $q \geqslant 0,01$: 
$\begin{aligned}[t]
B'(q) 
& = 8\times 2q\left (2-3\,\ln q\right ) +8q^2\left (0 -\dfrac{3}{q}\right ) - 0\\
& = 8q\left (4 - 6\,\ln q\right ) + 8q\left (-3\right )
= 8q \left (1 - 6\,\ln q\right )
\end{aligned}$
		
		\item 
$q>0$ donc $B'(q)$ est du signe de $1- 6\,\ln q$.

$1- 6\,\ln q>0
\iff
1> 6\,\ln q
\iff
\dfrac{1}{6} > \ln q
\iff
\e^{\frac{1}{6}} > q$

$\e^{\frac{1}{6}} \approx 1,2$ : $B\left (\e^{\frac{1}{6}} \right ) \approx 13,7$;
$B(0,01) \approx -3$

On dresse le tableau de variation complet de la fonction $B$.

\begin{center}
{\renewcommand{\arraystretch}{1.3}
\psset{nodesep=3pt,arrowsize=2pt 3}
\def\esp{\hspace*{2.5cm}}
\def\hauteur{0pt}
$\begin{array}{|c| l *4{c}|}
\hline
q & 0,01  & \esp & \e^{\frac{1}{6}}\approx 1,2 & \esp & +\infty \\ 
\hline
B'(q) &  &   \pmb{+} & \vline\hspace{-2.7pt}0 & \pmb{-} & \\ 
\hline
 & &  &   \Rnode{max}{\approx 13,7}  &  &   \\  
B(q) & &     &  &  &  \rule{0pt}{\hauteur} \\ 
 & \Rnode{min1}{\approx -3} &   &  &  &   \Rnode{min2}{-\infty} \rule{0pt}{\hauteur}    
 \ncline{->}{min1}{max} 
 \ncline{->}{max}{min2} 
 \\ 
\hline
\end{array} $
}
\end{center}	
		\item 
Le maximum de la fonction $B$ est $B\left (\e^{\frac{1}{6}} \right ) \approx 13,7$, donc le bénéfice maximum que peut espérer l'artisan est de 137~\euro.
	\end{enumerate}
\item 
	\begin{enumerate}
		\item 
\begin{list}{\textbullet}{Sur l'intervalle $[1,2~;~+\infty[$:}
\item La fonction $B$ est dérivable, donc continue, et strictement décroissante.		
\item Elle décroît de $13,7$ à $-\infty$.
\end{list}

Donc, d'après le corollaire des valeurs intermédiaires, l'équation $B(q)=10$ admet une solution unique sur l'intervalle $[1,2~;~+\infty[$. On l'appelle $\beta$.



$\left .
\begin{array}{l}
B(1) = 13 >10\\
B(2 \approx -5,5   <10
\end{array}
\right \}
\implies \beta \in [1 ~;~2]$
\hfill
$\left .
\begin{array}{l}
B(1,5) \approx 11,1 >10\\
B(1,6 \approx 9,1   <10
\end{array}
\right \}
\implies \beta \in [1,5 ~;~1,6]$

$\left .
\begin{array}{l}
B(1,55) \approx 10,17 >10\\
B(1,56 \approx 9,97   <10
\end{array}
\right \}
\implies \beta \in [1,55 ~;~1,56]$

$\left .
\begin{array}{l}
B(1,558) \approx 10,007 >10\\
B(1,559 \approx 9,986   <10
\end{array}
\right \}
\implies \beta \in [1,558 ~;~1,559]$

Donc $\beta$ a pour valeur approchée $1,558$ à $10^{-3}$ près.

		\item On admet que l'équation $B(q)= 10$ admet une unique solution $\alpha$ sur $[0,01~;~1,2[$.
		
On donne $\alpha \approx 0,757$, et on complète le tableau de variation de $B$.

\begin{center}
{\renewcommand{\arraystretch}{1.3}
\psset{nodesep=3pt,arrowsize=2pt 3}
\def\esp{\hspace*{2.5cm}}
\def\hauteur{20pt}
$\begin{array}{|c|*5{c}|}
\hline
q & 0,01  & \esp & 1,2 & \esp & +\infty \\ 
\hline
 & &  &   \Rnode{max}{13,7}  &  &   \\  
B(q) & &     &  &  &  \rule{0pt}{\hauteur} \\ 
 & \Rnode{min1}{-3} &   &  &  &   \Rnode{min2}{-\infty} \rule{0pt}{\hauteur}    
 \ncline{->}{min1}{max} 
 \ncline{->}{max}{min2} 
 \rput*(-6.1,0.95){\Rnode{zero}{\blue 10}}
\rput(-6.1,2.5){\Rnode{alpha}{\blue \alpha}}
\ncline[linestyle=dotted, linecolor=blue]{alpha}{zero}
\rput*(-2,0.95){\Rnode{zero2}{\red 10}}
\rput(-2,2.5){\Rnode{beta}{\red \beta}}
\ncline[linestyle=dotted, linecolor=red]{beta}{zero2}
 \\ 
\hline
\end{array} $
}
\end{center}	
		
Pour réaliser un bénéfice supérieur à $100$~euros, le nombre minimal de bonbons au chocolat à vendre correspond à $\alpha$ centaines, soit environ 76, et le nombre maximal de bonbons au chocolat à vendre correspond à $\beta$ centaines, soit environ 156.
	\end{enumerate}
\end{enumerate}
